\chapter{The Prakken--Sartor inference system: rules and fixed priorities}
\label{sec:p18-is1996}

Prakken and Sartor's \emph{A Dialectical Model of Assessing Conflicting
Arguments in Legal Reasoning}, in the held author version dated
July 19, 1996, constructs arguments from named rules, distinguishes
weak from strong negation, detects conflicts through strict extensions,
and compares the defeasible rules responsible for a conflict
\cite{PrakkenSartor1996Dialectical}.
Its technical companion supplies the fixed-point presentation and
the referred proof arguments \cite{PrakkenSartor1997Priorities}.
This chapter completes the fixed-priority construction and a corrected
finite proof procedure. The extension in which priorities are
themselves conclusions is supplied in
Chapter~\ref{sec:p18-is-dynamic}; its corrected full-domain
dialogue is proved in Chapter~\ref{sec:p18-is-dynamic-completeness}.

The source's arguments are rule sequences with potentially several
conclusions. They are not identified without explanation with the
single-rooted arguments of the separate ASPIC treatment.
The examples below are stipulated rule systems, including when
their atoms name familiar legal concepts; they state no current law.
General set-sized constructions use ordinary classical set theory
with choice; all choices in the finite examples are finite.

\section{Ground rules and arguments}

\begin{hbdefinition}{def:p18-is1996-theory}{Language and ordered theory}
Fix a first-order signature and its set of ground terms.
A \emph{strong literal} is a ground atom or its strong negation
\(\neg A\). Write \(\overline L\) for its complementary literal.
A \emph{weak literal} is \(\mathord{\sim}L\), expressing the
defeasible absence of support for \(L\), not strong support for
\(\overline L\).
A named rule has form
\[
r:L_1,\ldots,L_k,\mathord{\sim}M_1,\ldots,\mathord{\sim}M_j
   \Rightarrow N,
\]
where all displayed letters after the colon are strong literals.
Its name is a ground term. A variable-containing rule abbreviates
all its ground instances; use names containing the variables
needed to distinguish those instances.

An ordered theory is \(\Gamma=(S,\mathcal D,\prec)\).
The disjoint sets \(S,\mathcal D\) contain the strict and defeasible
ground rules. Strict rules contain no weak antecedent.
Write \(\to\) for membership in \(S\), as the source's notation
for this status, rather than treating it as an extra classical
implication connective. Empty antecedents are allowed.
Strict rules are not automatically closed under contraposition.
The relation \(\prec\) is a strict partial order on \(\mathcal D\);
\(r\prec s\) means that \(s\) has priority over \(r\).
No order on strict rules is required.
\end{hbdefinition}

\begin{hbdefinition}{def:p18-is1996-argument}{Arguments, conclusions, assumptions, and subarguments}
An argument is a finite list of rules such that every strong
antecedent of a rule was concluded by an earlier rule in the list,
and each conclusion has a unique producing rule in the list.
The empty argument \([]\) is allowed.
Weak antecedents are not proved when constructing the list.
They are recorded for later attack.

We use lists without redundant repetitions and identify valid
reorderings with the same rule content when comparing arguments.
This makes the source's content-based nonrepetition and its
set-inclusion comparisons explicit.
A subargument is a subset of the rule content that itself admits
a valid list, not an arbitrary subset lacking its required premises.
Write \(\mathcal A_\Gamma\) for the argument universe.

The set \(\operatorname{Con}(A)\) contains the consequents of
\emph{all} rules in \(A\).
The assumption set is
\[
\operatorname{Ass}(A)=
\{\overline M:\mathord{\sim}M\text{ occurs in a rule of }A\}.
\]
Thus a rule using \(\mathord{\sim}b\) assumes \(\neg b\);
it does not conclude \(\neg b\).
An argument is strict if all its rules are in \(S\), and
defeasible otherwise.
\end{hbdefinition}

\begin{hbexample}{def:p18-is1996-weak-strong}{Absence of evidence differs from contrary evidence}
Let
\[
r:\mathord{\sim}\operatorname{minor}(c)\Rightarrow
       \operatorname{capacity}(c),\qquad
s:\neg\operatorname{minor}(c)\Rightarrow\operatorname{vote}(c).
\]
With no further rules, \([r]\) is an argument and \([s]\) is not:
the latter lacks an earlier producer of its strong antecedent.
The conclusions of \([r]\) contain capacity alone.
Adding a rule \(\Rightarrow\operatorname{minor}(c)\) creates
an attacker of its weak assumption. Adding a rule
\(\Rightarrow\neg\operatorname{minor}(c)\) instead supplies the
missing premise for \(s\). These are different changes.
\end{hbexample}

\section{Strict extensions and conflicts}

\begin{hbdefinition}{def:p18-is1996-strict-closure}{Strict extensions and saturation}
A strict extension \(A+T\) appends strict rules so that the result
is an argument. Its added rules must retain the unique-producer
condition. Let
\[
C_S(A)=\{L:\text{some strict extension of }A
                    \text{ has conclusion }L\}.
\]
Equivalently, start with \(\operatorname{Con}(A)\) and repeatedly
apply ground strict rules whose strong premises have been obtained.
For infinite \(S\), this is a definition by finite derivations,
not a claimed terminating computation.
\end{hbdefinition}

\begin{hbtheorem}{thm:p18-is1996-closure}{proved=annex:p18.is1996-closure}{Saturation and extension agree}
The two descriptions of \(C_S(A)\) agree. If \(A\) is a
subargument of \(B\), then \(C_S(A)\subseteq C_S(B)\) and
\(\operatorname{Ass}(A)\subseteq\operatorname{Ass}(B)\).
Removing repeated occurrences of an identical producing rule
does not change the conclusions or assumptions of a valid list.
\end{hbtheorem}

\begin{hbdefinition}{def:p18-is1996-attack}{Head conflict, assumption attack, and coherence}
Argument \(A\) attacks \(B\) if either:
\begin{enumerate}
\item there is \(L\in C_S(A)\) with
\(\overline L\in C_S(B)\); or
\item there is \(L\in C_S(A)\) with
\(\overline L\in\operatorname{Ass}(B)\).
\end{enumerate}
The first is a head conflict and is symmetric.
The second is the source's \emph{undercutting} relation
\(U(A,B)\): a weak antecedent is contradicted.
It need not be symmetric, although mutual undercutting is possible.
This source-specific usage need not coincide with later uses of
the word undercutter.
An argument is coherent when it does not attack itself.
A set is conflict-free when no member attacks any member,
including itself.
\end{hbdefinition}

\begin{hbexample}{def:p18-is1996-strict-conflict}{A conflict exposed by a strict convention}
Let \(r:\Rightarrow\operatorname{lease}(c)\) and
\(s:\Rightarrow\neg\operatorname{contract}(c)\) be defeasible,
and let
\(\operatorname{lease}(x)\to\operatorname{contract}(x)\) be strict.
The arguments \([r]\) and \([s]\) attack each other even though
their immediate consequents are not complementary.
The strict extension of \([r]\) exposes the conflict.
The priority comparison must concern \(r\) and \(s\);
it must not demote the strict convention to a defeasible rival.
\end{hbexample}

\section{Relevant rules and defeat}

\begin{hbdefinition}{def:p18-is1996-frontier}{The relevant defeasible frontier}
For a fixed argument \(E\) and \(L\in\operatorname{Con}(E)\),
define \(R_L(E)\) recursively along its unique producing rule.
If the producer is defeasible \(r\), set \(R_L(E)=\{r\}\).
If the producer is strict, with strong antecedents
\(L_1,\ldots,L_k\), set
\[
R_L(E)=\bigcup_{i=1}^k R_{L_i}(E).
\]
For a strict empty-premise rule the union is empty.
Apply this recurrence to the whole strictly extended argument,
including strict rules already present in its original prefix.
The frontier is unique for that argument; different strict
extensions can have different frontiers and must not be silently merged.
\end{hbdefinition}

\begin{hbtheorem}{thm:p18-is1996-frontier}{proved=annex:p18.is1996-frontier}{Frontier construction terminates}
The recurrence determines a unique finite set of defeasible rules
for every produced literal. A strict argument has empty frontiers.
A defeasible final producer contributes itself, regardless of
earlier defeasible support; attacks on that earlier support
remain possible through the other conclusions of the argument.
\end{hbtheorem}

\begin{hbdefinition}{def:p18-is1996-set-priority}{The source's strict set comparison}
For finite sets \(R,R'\subseteq\mathcal D\), define
\[
R\prec_*R'
\quad\Longleftrightarrow\quad
\exists r\in R\ \forall r'\in R'\ (r\prec r').
\]
One rule on the left must be below every rule on the right.
This is not a comparison of average weights or of all pairs.
\end{hbdefinition}

\begin{hbtheorem}{thm:p18-is1996-set-priority}{proved=annex:p18.is1996-set-priority}{Order and empty-frontier cases}
The lifted relation is a strict partial order.
Every nonempty \(R\) satisfies \(R\prec_*\varnothing\);
the empty set is below no set.
In particular, both \(R\prec_*R'\) and \(R'\prec_*R\)
cannot hold.
\end{hbtheorem}

\begin{hbdefinition}{def:p18-is1996-defeat}{Rebuttal and the generated defeat relation}
Argument \(A\) rebuts \(B\), written \(R(A,B)\), when there
are strict extensions \(A+T,B+V\) and complementary conclusions
\(L,\overline L\) for which
\[
R_L(A+T)\not\prec_*R_{\overline L}(B+V).
\]
The quantifier is existential over conflict witnesses and their
actual strict derivations.
Define \(D_\Gamma(A,B)\), or simply \(A D_\Gamma B\), by:
\[
\bigl(A=[]\text{ and }B\text{ incoherent}\bigr)
\quad\text{or}\quad U(A,B)
\quad\text{or}\quad\bigl(R(A,B)\text{ and not }U(B,A)\bigr).
\]
Strict defeat \(A D_\Gamma^s B\) means
\(A D_\Gamma B\) and not \(B D_\Gamma A\).
The first disjunct is a designated incoherence defeat.
It need not itself be an attack by the empty argument.
\end{hbdefinition}

\begin{hbtheorem}{thm:p18-is1996-orientation}{proved=annex:p18.is1996-orientation}{Conflict orientation and inherited attacks}
Every attack between two arguments yields a defeat in at least
one direction. Rebuttal and undercutting imply attack.
Defeat implies attack except possibly for the designated empty
incoherence case.
An attack on a subargument is an attack on its superargument,
although its direction of successful defeat may change.
If strict arguments are not conflict-free, every argument is incoherent.
\end{hbtheorem}

\begin{hbexample}{def:p18-is1996-undercut}{Undercutting takes precedence over rebuttal}
Take \(r:\mathord{\sim}\neg a\Rightarrow a\) and
\(s:\Rightarrow\neg a\), with no priorities.
Each argument rebuts the other. But \([s]\) undercuts \([r]\),
whereas \([r]\) does not undercut \([s]\).
Thus \([s]\) strictly defeats \([r]\).
By contrast, \([\mathord{\sim}b\Rightarrow a]\) and
\([\mathord{\sim}a\Rightarrow b]\) mutually undercut;
neither strictly defeats the other.
\end{hbexample}

\section{The original fixed-point semantics}
\label{sec:p18-is1996-semantics}

For \(X\subseteq\mathcal A_\Gamma\), put
\[
F_\Gamma(X)=\{A:\ \forall B\,
 [B D_\Gamma A\Rightarrow
       \exists C\in X\ (C D_\Gamma^s B)]\}.
\]
The source requires strict defense against every defeater.
Let \(J_\Gamma\) be its least fixed point, obtained by iteration
from the empty set, using unions at limit ordinals.
An ordered theory is \emph{finitary} here when every argument
has only finitely many attackers. The extra empty incoherence
defeater adds at most one further defeater.

\begin{hbtheorem}{thm:p18-is1996-semantics}{proved=annex:p18.is1996-semantics}{Existence, finite approximation, and structural properties}
\(F_\Gamma\) is monotone and has a unique least fixed point.
Its finite iterates are contained in \(J_\Gamma\).
For finitary \(\Gamma\), their union equals \(J_\Gamma\).
The set \(J_\Gamma\) is conflict-free under the source's attack
relation and contains every subargument of each of its members.
If the strict arguments are not conflict-free, \(J_\Gamma=\varnothing\).
\end{hbtheorem}

The source calls members of \(J_\Gamma\) justified; it calls an
argument overruled when it is attacked by a justified argument,
and defensible when it is neither justified nor overruled.
Retain the word \emph{attacked} in this source-specific definition.
Because the empty incoherence defeat need not be an attack,
the residual class can contain incoherent arguments.
It is a technical residual classification, not a claim that
such an argument is a reasonable position.
If the strict base is inconsistent, every argument falls in
that residual class because there are no justified arguments.

A literal is justified if some justified argument concludes it;
it is defensible if it is not justified and some defensible
argument concludes it; and it is overruled if neither prior
case applies but an overruled argument concludes it.
If no argument concludes it, record it as unsupported.
These classifications do not assert empirical truth or falsity.

\section{A corrected finite proof game}

\begin{hbdefinition}{def:p18-is1996-proof-game}{Finite dialectical proofs}
A dialogue alternates proponent \(P\) and opponent \(O\),
starting with the argument to be proved.
Every subsequent \(O\) move is any \(D_\Gamma\)-defeater
of the preceding argument. Every subsequent \(P\) move is
a \(D_\Gamma^s\)-defeater of the preceding argument.
\(P\) may not repeat the same argument content on one branch;
\(O\) may repeat.
There is no additional global minimality restriction on \(P\)'s reply.

A finite proof tree has one such dialogue on each branch,
contains every possible opponent reply at each \(P\) node,
and chooses a successful \(P\) response at each nonterminal
\(O\) node. It is won by \(P\) when every leaf is a \(P\) move
with no opponent reply. A terminal play alone is not a proof
unless all opponent alternatives have been covered.
\end{hbdefinition}

\begin{hbtheorem}{thm:p18-is1996-proof-correctness}{proved=annex:p18.is1996-proof-correctness}{General soundness and finitary completeness}
Every argument with a finite proof tree in this corrected game
belongs to \(J_\Gamma\), for arbitrary \(\Gamma\).
For finitary \(\Gamma\), every member of \(J_\Gamma\) has
such a finite proof tree.
\end{hbtheorem}

\begin{hbtheorem}{thm:p18-is1996-minimality}{proved=annex:p18.is1996-minimality}{Global reply minimality can destroy completeness}
If \(P\)'s replies must instead be globally minimal, by inclusion
of rule content, among strict defeaters of the preceding argument,
finitary completeness and proof-theoretic subargument closure fail.
A four-rule finite theory with no strict rules or priorities
provides a counterexample.
\end{hbtheorem}

\begin{hbexample}{def:p18-is1996-minimality-case}{Two pieces of support can protect each other}
Use exactly
\[
r_1:\Rightarrow a,\quad r_2:\Rightarrow b,\quad
r_3:\mathord{\sim}b\Rightarrow\neg a,\quad
r_4:\mathord{\sim}a\Rightarrow\neg b.
\]
The combined argument \(C=[r_1,r_2]\) has no defeaters:
every argument with a contrary conclusion uses \(r_3\) or \(r_4\),
whose weak antecedent \(C\) undercuts.
Hence \(C\) is justified immediately and its subarguments
\([r_1]\), \([r_2]\) are justified as well.

With global minimality, an opponent can force
\[
P:[r_1],\quad O:[r_3],\quad P:[r_2],\quad O:[r_4].
\]
The only minimal strict response now repeats \([r_1]\).
The protective \(C\) is not globally minimal, so that version
cannot finish a proof of \([r_1]\).
In the corrected game, \(P\) replies with \(C\) and wins.
The annex proves the claim against all possible argument contents;
it is not a comparison of an arbitrarily selected defeat graph.
\end{hbexample}

\begin{hbexample}{def:p18-is1996-earlier-support}{An earlier defeat can settle a later conflict}
Take
\[
r_1:\Rightarrow a,\quad r_2:a\Rightarrow b,\quad r_3:b\Rightarrow c,
\qquad
r_4:\Rightarrow\neg a,\quad r_5:\neg a\Rightarrow d,\quad
r_6:d\Rightarrow\neg c
\]
with \(r_4\prec r_1\) and \(r_3\prec r_6\).
The complete chains \([r_1,r_2,r_3]\) and \([r_4,r_5,r_6]\)
can defeat each other through different conclusions.
The single argument \([r_1]\), however, strictly defeats every
argument containing \(r_4\), and has no defeaters.
Every attack against the positive chain requires \(r_4\):
the contrary conclusions \(\neg a,\neg c\) both arise through it.
Thus respond to each such opponent with \([r_1]\).
This gives a complete finite proof, covering all opposing
superarguments as well as the displayed negative chain.
\end{hbexample}

\section{Finite computation and the LIS connection}

\begin{hbtheorem}{thm:p18-is1996-finite}{proved=annex:p18.is1996-finite}{A decision procedure for a finite ground theory}
If \(S\cup\mathcal D\) is a finite set of ground rules and
its priority relation is given, the argument universe, attack
and defeat relations, fixed-point statuses, and corrected
finite proof outcomes can all be computed by finite enumeration.
Finitely many written schemata with function symbols do not
by themselves satisfy this hypothesis.
\end{hbtheorem}

For example, \(\to p(c)\) and \(p(x)\Rightarrow p(f(x))\)
are only two written schemata but produce arbitrarily long
arguments with different conclusions \(p(f^n(c))\).
Conversely, a theory can have infinitely many arguments and
no attacks, making it finitary without being finite.
The finite-ground algorithm must not be described as an
algorithm for every finitary theory.

\begin{hbtheorem}{thm:p18-is1996-lis-bridge}{proved=annex:p18.is1996-lis-bridge}{The precise fixed-role bridge}
Use the actual rule-generated relation \(D_\Gamma\) as the
defeat graph for the corrected fixed-role game of
Chapter~\ref{sec:p18-lis-burdens}.
Its grounded extension equals \(J_\Gamma\).
For finitary \(\Gamma\), its winning arguments are exactly
those with the corrected finite proofs above.
This identifies justification only; it does not identify the
source's attack-based residual labels with every later label convention.
\end{hbtheorem}

There is no identity bridge to the changing-role LIS interface
without extra data. A rule-sequence argument can have several
conclusions and a counterargument can attack several incomparable
subarguments. For instance, separate conclusions \(a,b\)
can be opposed by an argument containing both \(\neg a,\neg b\).
The single-target role rule in that LIS interface does not
choose a target for this case. A designated conclusion, occurrence
representation, and justified target policy are additional requirements.

\begin{proofboundary}{def:p18-is1996-next}{Exact continuing source obligations}
This bounded component supplies the original ground-rule construction,
strict-extension conflicts, relevant frontiers, fixed priorities,
generated defeat, fixed-point semantics, and a corrected finite
proof theorem with the original general/finitary scope.
The global-minimality formulation is refuted by the full
rule-generated counterexample rather than asserted complete.

Chapter~\ref{sec:p18-is-dynamic} supplies the reversed-role
fixed-comparison game, derived priorities, neutralization,
the dynamic semantic proofs, the priority examples and the
applicability alternatives with their exact failure boundaries.
The original dynamic game's finite completeness claim is
refuted by Theorem~\ref{thm:p18-is-compression-countermodel},
so an unrestricted finitary extension of that claim is also false.
The corrected sequential game is complete on the original
set-sized domain by Theorem~\ref{thm:p18-is-licensing-complete},
with finite proofs for finitary theories by
Theorem~\ref{thm:p18-is-licensing-finite}.

The separate sensitive-exclusion semantics, unrestricted
non-semidecidability coding, broader stable-semantics comparisons
and historical Dung-version question retain their stated
source obligations.
\end{proofboundary}
